\section*{Chapter \ref{ch:g9:institutions-and-prosperity} --- Why Some Countries Are Rich: Institutions}

\begin{solution}{exo:g9:institutions-and-prosperity:1}
In 1940 the north (2,624 dollars against 1,975). In 2022 the south (41,321 against 1,569).
\end{solution}

\begin{solution}{exo:g9:institutions-and-prosperity:2}
$41\,321 \div 1\,569 = 26.3$: about 26 times as much.
\end{solution}

\begin{solution}{exo:g9:institutions-and-prosperity:3}
$19\,238 \div 1\,744 = 11.03$: about 11 times.
\end{solution}

\begin{solution}{exo:g9:institutions-and-prosperity:4}
$11.03$ lies between $10.55$ ($5.5\,\%$) and $12.98$ ($6\,\%$), closer to the first: China grew by a little
more than $5.5\,\%$ a year (more precisely, $5.6\,\%$).
\end{solution}

\begin{solution}{exo:g9:institutions-and-prosperity:5}
South: $41\,321 \div 13\,874 = 2.98$, close to $3.01$: about $3.5\,\%$ a year. North: $1\,569 \div 2\,456 =
0.64$: a fall of about $1.4\,\%$ a year.
\end{solution}

\begin{solution}{exo:g9:institutions-and-prosperity:6}
In the poorest countries, often between 10 and $30\,\%$ of \omterm{def:g7:firms:firm}{firms}, with a middle value of $14.4\,\%$ for the
countries under 5,000 dollars per person; in the richest, a few percent or less, with a middle value of
$2.6\,\%$ above 30,000 dollars.
\end{solution}

\begin{solution}{exo:g9:institutions-and-prosperity:7}
The model predicts more farm output: each extra sack grown beyond the delivery is the family's to keep or
sell. In China, the return of farms to families accounted for about half of the growth of farm output from
1978 to 1984 (Lin, 1992).
\end{solution}

\begin{solution}{exo:g9:institutions-and-prosperity:8}
1991: $100 \div 40 = 2.5$ times as much. 1996: $100 \div 60 = 1.67$ times, or a little less, since the share
was somewhat above $60\,\%$.
\end{solution}

\begin{solution}{exo:g9:institutions-and-prosperity:9}
For example: two parts of a country with the same people and land, one under a \omterm{def:g9:institutions-and-prosperity:planned-economy}{planned economy} and one under
a \omterm{def:g9:institutions-and-prosperity:market-economy}{market economy}, where the planned part grows faster for decades; or a reform that secures farmers'
property and lets \omterm{def:g7:money:price}{prices} move, followed by lower farm output for many years with nothing else changing.
\end{solution}

\begin{solution}{exo:g9:institutions-and-prosperity:10}
The table cannot tell. Courts and \omterm{def:g7:households:income}{income} move together, but so do the coast and \omterm{def:g7:households:income}{income}: access to the sea
makes trade cheaper and could make the coastal countries rich, whatever their courts. Riches could also pay
for better courts. To separate the causes one needs countries alike in everything but their courts.
\end{solution}

\begin{solution}{exo:g9:institutions-and-prosperity:11}
Lower farm output, or slower growth, over the following decades: farmers no longer keep what extra effort or
investment produces, and fixed \omterm{def:g7:money:price}{prices} no longer tell them what is most wanted, so they work, invest and
adapt less.
\end{solution}

\begin{solution}{exo:g9:institutions-and-prosperity:12}
It shows a \omterm{def:g9:institutions-and-prosperity:correlation}{correlation} only. Poverty could cause bribery (badly paid officials ask for bribes), or a third
cause, such as war or weak schooling, could cause both poverty and bribery.
\end{solution}

\begin{solution}{pb:g9:institutions-and-prosperity:1}
\textbf{1.} All 100 dalens: the rent is fixed.
\textbf{2.} $100 \div 100 = 1$ dalen: the extra harvest is shared by all 100 families.
\textbf{3.} Valley A: each family keeps what its extra effort produces, so families work and invest more.
\textbf{4.} A: $1\,000 \times 1.03 = 1\,030$, then $1\,030 \times 1.03 = 1\,061$. B: $1\,000 \times 1.01 =
1\,010$, then $1\,010 \times 1.01 = 1\,020$.
\textbf{5.} $1.03^{20} = 1.806$, so $1\,000 \times 1.806 = 1\,806$.
\textbf{6.} $1\,806 \div 1\,220 = 1.48$: $1.48$ times as much, $48\,\%$ ahead.
\textbf{7.} Year 21: A $1\,806 \times 1.03 = 1\,860$, B $1\,220 \times 1.05 = 1\,281$. Year 22: A
$1\,860 \times 1.03 = 1\,916$, B $1\,281 \times 1.05 = 1\,345$.
\textbf{8.} $1\,916 \div 1\,345 = 1.425$: $42.5\,\%$ ahead.
\textbf{9.} $1.03 \div 1.05 = 0.981$: the ratio shrinks by about $2\,\%$ a year.
\textbf{10.} $24\,\%$.
\textbf{11.} Year 29: $2\,357 \div 1\,893 = 1.245$, a lead of $24.5\,\%$, still above $24\,\%$. Year 30:
$2\,427 \div 1\,988 = 1.221$, a lead of $22.1\,\%$. The lead is halved for the first time in year 30.
\textbf{12.} Ten years after the reform. In dalens the gap shrinks more slowly (586 in year 20, $2\,427 -
1\,988 = 439$ in year 30), because valley A keeps growing by $3\,\%$ on a larger amount: its yearly gain in
dalens stays large even while B grows faster in percent.
\end{solution}
